\section{Escribir KalaScript en LaTeX}

Uso los comandos de esta sección como herramientas de edición. Sus nombres no
introducen pronunciaciones ni forman parte de la escritura manuscrita.

\subsection{Un signo y una referencia}

\begin{Verbatim}[fontsize=\small,frame=single]
\KalaSimbolo{K01}
\KalaSimbolo[15mm]{K01}
\KalaReferencia{1}
\end{Verbatim}

\KalaMostrar{\KalaSimbolo{K01}\quad\KalaSimbolo[15mm]{K01}\quad\KalaReferencia{1}}

El tamaño se puede ajustar sin cambiar la proporción del trazo. La referencia
necesita una definición en la leyenda del documento.

\subsection{Componer y delimitar}

\begin{Verbatim}[fontsize=\small,frame=single]
\KalaNegacion{
  \KalaConjuncion{\KalaMat{P}}{\KalaMat{Q}}
}
\end{Verbatim}

\KalaMostrar{\KalaNegacion{\KalaConjuncion{\KalaMat{P}}{\KalaMat{Q}}}}

\texttt{\textbackslash KalaMat} recibe una expresión matemática convencional,
sin delimitadores de modo matemático adicionales. Los conectivos reciben
expresiones completas como argumentos.

\subsection{Cuantificar}

\begin{Verbatim}[fontsize=\small,frame=single]
\KalaUniversal{x}{\mathbb{N}}{
  \KalaExistencia{y}{\mathbb{N}}{
    \KalaMat{x < y}
  }
}
\end{Verbatim}

\KalaMostrar{\KalaUniversal{x}{\mathbb{N}}{\KalaExistencia{y}{\mathbb{N}}{\KalaMat{x<y}}}}

\subsection{Documentar una relación}

\begin{Verbatim}[fontsize=\small,frame=single]
\KalaRelacion{2}{
  \KalaReferencia{1}\KalaSeparador
  \KalaReferencia{2}\KalaSeparador
  \KalaReferencia{3}
}
\end{Verbatim}

Los índices, los referentes y el orden de argumentos deben corresponder al
catálogo y a la leyenda. La compilación comprueba que los comandos existan y
que LaTeX pueda representarlos; no comprueba la verdad, el dominio, las
unidades ni la validez semántica de cada relación.

\subsection{Modificar un signo}

En \path{simbolos/glifos.tex}, cada llamada a
\texttt{\textbackslash KalaRegistrar} contiene seis argumentos:
código, nombre, familia, definición, número de trayectos y trazado TikZ.
Las coordenadas siguen la cuadrícula de 64 por 64 unidades. Las curvas
cuadráticas originales se expresan como curvas cúbicas equivalentes.

Al modificar una grafía, recompilo los cinco documentos y reviso las
expresiones donde aparece. Mantengo el identificador si la identidad del signo
no cambia y documento por separado cualquier modificación de significado.

\subsection{Composición de expresiones}
\begin{Verbatim}[fontsize=\small,frame=single]
\KalaConcepto{3}     % sensores del estudio
\KalaPropiedad{2}    % activo, contexto declarado
\KalaSeleccion{x}{\KalaConcepto{3}}{
  \KalaPosee{\KalaVariable{x}}{2}
}
\KalaAlternativa{\KalaMat{P}}{\KalaMat{Q}}
\KalaPosible{1}{\KalaMat{P}}
\KalaPregunta{\KalaMat{P}}
\KalaSecuencia{\KalaAsignar{x}{0}}{
  \KalaMientras{\KalaMat{x<3}}{\KalaAsignar{x}{x+1}}
}
\end{Verbatim}
\texttt{\textbackslash KalaTodo} y \texttt{\textbackslash KalaExiste} reciben
una variable, un dominio ya dibujado y un cuerpo. Los comandos anteriores
\texttt{\textbackslash KalaUniversal} y \texttt{\textbackslash KalaExistencia}
siguen recibiendo el dominio como matemática convencional sin dibujar.

\texttt{\textbackslash KalaColeccion} recibe tipo (1 conjunto, 2 lista) y miembros
separados con \texttt{\textbackslash KalaSeparador}. Los alias gráficos y comandos de formas anteriores se documentan en el anexo
de compatibilidad; para textos nuevos uso los comandos habituales.

Los trazados K17--K22 están en \path{simbolos/extension-02.tex} y sus comandos en
\path{comandos-02.tex}. El banco \path{datos/casos-02.tex} contiene para cada caso:
identificador, área, intención, expresión, contexto, lectura, contraste y cobertura.
